\documentclass[11pt,a4paper]{article}

\usepackage[utf8]{inputenc}
\usepackage[T1]{fontenc}
\usepackage[brazil]{babel}
\usepackage{lmodern}
\usepackage[margin=2.3cm]{geometry}
\usepackage{booktabs}
\usepackage{tabularx}
\usepackage{array}
\usepackage{xcolor}
\usepackage{colortbl}
\usepackage{enumitem}
\usepackage{hyperref}
\usepackage{caption}
\usepackage{tikz}
\usepackage{pgfplots}
\usepackage{siunitx}
\usepackage{float}
\usepackage{xurl}

\pgfplotsset{compat=1.18}
\usetikzlibrary{positioning,arrows.meta}
\sisetup{output-decimal-marker={,}, group-separator={.}, group-minimum-digits=4}

\definecolor{cblue}{HTML}{2A78D6}
\definecolor{corange}{HTML}{EB6834}
\definecolor{caqua}{HTML}{1BAF7A}
\definecolor{cgray}{HTML}{8A8984}
\definecolor{rlow}{HTML}{D9F2E4}
\definecolor{rmed}{HTML}{FCEBC4}
\definecolor{rhigh}{HTML}{F8D3CF}
\definecolor{boxbg}{HTML}{EEF4FC}

\hypersetup{colorlinks=true, linkcolor=cblue, urlcolor=cblue}
\setlist{nosep, leftmargin=1.4em}
\captionsetup{font=small, labelfont=bf}
\newcolumntype{L}{>{\raggedright\arraybackslash}X}
\renewcommand{\arraystretch}{1.15}

\newcommand{\Rb}{\cellcolor{rlow}Baixo}
\newcommand{\Rm}{\cellcolor{rmed}Médio}
\newcommand{\Ra}{\cellcolor{rhigh}Alto}

\title{\textbf{Previsão do tempo de entrega (Meituan)}\\[2pt]
\large Relatório final — Etapa 4: Governança da Solução e MLOps}
\author{Guilherme Braga \and Guilherme Costanzi \and Rafael Cobo\\[2pt]
\small Governança de TI e Dados — Universidade Estadual de Londrina (UEL)}
\date{Novembro de 2026}

\begin{document}
\maketitle

\begin{abstract}
\noindent Este relatório reúne o pacote de governança da nossa solução de previsão do tempo de
entrega, feita com os dados públicos da Meituan. Ele junta o que entregamos nas etapas
anteriores (documentação dos dados e pipeline com DVC e MLflow) com os resultados da segunda
rodada de testes de modelos, o model card, o plano de monitoramento, a matriz de riscos e as
regras de governança da solução. A Seção~\ref{sec:checklist} mostra onde está cada item pedido
na entrega final.
\end{abstract}

%=====================================================================
\section{Relatório executivo}
\label{sec:exec}

\paragraph{Problema.} Num aplicativo de delivery, o cliente recebe uma estimativa de quando o
pedido vai chegar. Se a estimativa de tempo é pequena demais, o pedido pode atrasar gerar insatisfação ao cliente. Se é
longa demais, o serviço parece lento e o cliente pode desistir da compra. Na base da Meituan, a estimativa atual
do aplicativo erra em média \textbf{9,5 minutos}. Isso acontece propositalmente para evitar atrasos
e gerar folga na previsão. O pedido chega, em média, 8 minutos antes do prometido.

\paragraph{Objetivo.} Prever, em minutos, quanto tempo passa entre a criação do pedido e a
chegada ao cliente. A previsão é feita em dois momentos: na \textbf{compra} (\emph{checkout},
quando ainda não há entregador) e no \textbf{despacho} (quando o entregador aceita o pedido).
Além de obter uma previsão precisa, temos como objetivo conseguir explicar, repetir e controlar
como o modelo foi feito, garantindo governança e rastreabilidade ao trabalho desenvolvido.

\paragraph{Resultado.} Nossos modelos erram, em média, cerca de \textbf{5 minutos}, quase
metade do erro da estimativa atual (Tabela~\ref{tab:r2}). O modelo que menos errou foi o TabPFN
(5,006 min). O LightGBM teve resultado próximo (5,133 min), com custo computacional muito menor.
As transformações nos dados foram documentadas com a ferramenta \emph{DVC} e a rastreabilidade dos
modelos foi garantida com a ferramenta \emph{MLflow} em um servidor privado. 

\paragraph{Decisão.} Indicamos o \textbf{LightGBM} como o modelo para uso (\emph{candidato}),
o \textbf{TabPFN} como \emph{desafiante} (o modelo que pode substituí-lo no futuro, se provar
que é melhor) e o \textbf{Random Forest} como \emph{referência}, por ter sido o modelo da
demonstração do MLflow. O modelo serve para apoiar decisões. Ele não define sozinho o prazo
mostrado ao cliente e também não pode ser usado para avaliar entregadores.

\paragraph{Como a solução é governada.} O código fica no Git; os dados e o passo a passo do
processamento, no DVC; os testes de modelos, no MLflow. O model card explica o que o modelo
faz e o que não faz. Nosso plano de monitoramento diz o que medir, qual os limites que engatilham ações
e o que fazer quando esses gatilhos são atingidos.
Há uma matriz de riscos que lista os riscos e define responsáveis por cuidar de cada um. 
Nenhuma versão deve entrar em uso sem aprovação registrada (Seção~\ref{sec:gov}).

\paragraph{Ligação com os objetivos estratégicos.}
\begin{itemize}
  \item \textbf{Experiência do cliente:} prazos mais próximos do real reduzem atrasos e
        reclamações.
  \item \textbf{Eficiência da operação:} saber o tempo esperado ajuda a planejar quantos
        entregadores são necessários em cada região e horário.
  \item \textbf{Confiança e conformidade:} saber de onde vem cada resultado e ter regras de uso
        claras reduz riscos legais (licença dos dados, privacidade da localização) e éticos.
\end{itemize}

%=====================================================================
\section{Repositório e reprodução}
\label{sec:repo}

O código começou num notebook. Para que qualquer pessoa consiga repetir o trabalho, ele foi
separado em um pacote Python (\texttt{meituan\_eta}) e cada passo virou uma etapa do DVC
(Figura~\ref{fig:pipe}). O DVC sabe o que cada etapa recebe como entrada e o que gera de saída,
criando um grafo acíclico de dependências entre os dados e scripts. Se algo alguma dependência
muda, a ferramenta consegue refazer apenas as etapas necessárias.

\begin{figure}[H]
\centering
\begin{tikzpicture}[node distance=0.35cm,
  box/.style={draw=cblue, rounded corners=3pt, minimum width=1.85cm, minimum height=0.9cm,
              font=\footnotesize, align=center, line width=0.7pt},
  arr/.style={-{Stealth[length=2mm]}, cgray, line width=0.7pt}]
\node[box] (b) {bootstrap\\\scriptsize dados brutos};
\node[box, right=of b] (v) {validate\\\scriptsize validação};
\node[box, right=of v] (p) {prepare\\\scriptsize limpeza};
\node[box, right=of p] (f) {features\\\scriptsize variáveis};
\node[box, right=of f] (s) {split\\\scriptsize divisão};
\node[box, right=of s] (t) {train\\\scriptsize treino};
\node[box, right=of t] (e) {evaluate\\\scriptsize avaliação};
\foreach \a/\b in {b/v,v/p,p/f,f/s,s/t,t/e} \draw[arr] (\a) -- (\b);
\end{tikzpicture}
\caption{Etapas do pipeline no DVC (arquivo \texttt{dvc.yaml}).}
\label{fig:pipe}
\end{figure}

\noindent Para repetir o trabalho, a partir da pasta principal do repositório:
\begin{enumerate}
  \item criar o ambiente Python 3.12 e instalar as bibliotecas de \texttt{requirements.txt}
        (com versões fixas);
  \item colocar os quatro arquivos originais do desafio \emph{Meituan INFORMS TSL Research Challenge}
  em \texttt{data/} (eles não foram integrados ao repositório porque a licença não permite);
  \item rodar \texttt{dvc repro} para executar o pipeline e \texttt{dvc metrics show} para ver
        os resultados;
  \item abrir o MLflow com \texttt{mlflow server --backend-store-uri sqlite:///mlflow.db}.
\end{enumerate}
Os testes da rodada 2 estão nos notebooks \texttt{notebooks/5\_meituan\_tabpfn\_colab*.ipynb},
que usam \texttt{notebooks/dataloader.py}. O \texttt{README.md} explica cada comando, e a pasta
\texttt{tests/} tem testes automáticos do código.

\subsection{Versionamento}

Versionar é guardar cada versão de algo de forma que dê para voltar a ela e saber o que mudou.
A Tabela~\ref{tab:versao} mostra como fazemos isso com cada parte da solução.

\begin{table}[H]
\centering\small
\begin{tabularx}{\linewidth}{l l L}
\toprule
O que & Ferramenta & Como identificamos a versão \\
\midrule
Código & Git & cada mudança vira um \emph{commit}; o MLflow guarda o commit usado em cada execução. \\
Dados & DVC & o \texttt{dvc.lock} guarda o hash de cada arquivo de dados; o MLflow guarda o hash do conjunto usado no treinamento do modelo. \\
Configurações & Git + DVC & \texttt{params.yaml} guarda parâmetros, semente aleatória e a divisão dos dados; se um valor mudar, o DVC pode refazer a etapa. \\
Modelo & DVC + MLflow & o arquivo do modelo (\texttt{models/*.pkl}) tem hash no \texttt{dvc.lock} e uma cópia na execução do MLflow. \\
Ambiente & Git & \texttt{requirements.txt} com as versões das bibliotecas. \\
\bottomrule
\end{tabularx}
\caption{Formas de versionamento para cada domínio da solução implementada.}
\label{tab:versao}
\end{table}

%=====================================================================
\section{Dataset documentado}
\label{sec:dados}

\paragraph{Origem.} Os dados vêm do \emph{INFORMS TSL Data-Driven Research Challenge} e foram
fornecidos pela Meituan, uma empresa chinesa de delivery.

\begin{table}[H]
\centering\small
\begin{tabularx}{\linewidth}{l L}
\toprule
Período & 17 a 24 de outubro de 2022 (8 dias), uma cidade, horário da China (UTC+8) \\
Volume & 654.343 registros de entrega com 23 colunas, além de tabelas de ondas de entrega e de despacho \\
O que é uma linha & uma tentativa de entrega; se um entregador recusa o pedido, este pedido pode aparecer de novo \\
Privacidade & os IDs foram trocados por códigos; as coordenadas foram deslocadas, mas não embaralhadas \\
Licença & CC BY-NC 4.0: só uso não comercial não sendo permitido redistribuir os dados originais \\
\bottomrule
\end{tabularx}
\caption{Origem dos dados. O dicionário de dados está no Apêndice~\ref{sec:dic}.}
\end{table}

\paragraph{Qualidade.} Não há células vazias, mas alguns campos usam o número zero para dizer
``sem informação'' (1,51\% no tempo estimado de preparo). Não há linhas duplicadas, os
horários estão na ordem certa e não há velocidades impossíveis. Os detalhes estão em
\texttt{quality/QUALITY\_REPORT.md}.

\paragraph{Transformações.} A Tabela~\ref{tab:funil} mostra quantos registros sobram após cada
filtro. O valor que queremos prever é o tempo, em minutos, entre a criação do pedido e a
chegada ao cliente.

\begin{table}[H]
\centering\small
\begin{tabular}{l r r}
\toprule
Filtro & Registros que ficam & Removidos \\
\midrule
Base original & \num{654343} & -- \\
Só pedidos aceitos por um entregador & \num{568546} & \num{85797} \\
Só pedidos imediatos (sem agendamento) & \num{546365} & \num{22181} \\
Só horários válidos (sem zeros) & \num{538920} & \num{7445} \\
Só entregas entre 5 e 90 minutos & \textbf{\num{538435}} & 485 \\
\bottomrule
\end{tabular}
\caption{Limpeza dos dados: ficam 82,3\% dos registros.}
\label{tab:funil}
\end{table}

\paragraph{Variáveis criadas.} A partir das colunas originais, o carregador de dados
do projeto (arquivo \texttt{dataloader.py}, na pasta \texttt{notebooks}) cria três conjuntos de variáveis: 13 disponíveis no momento
da compra (distâncias pela métrica de Manhattan, horário do pedido, tempo estimado de preparo,
distrito e tempo médio histórico do restaurante); 25 no momento do despacho (acrescenta posição
e carga do entregador, tempos até o aceite e recusas anteriores); e 38 na rodada 2, com
variáveis sobre a situação do restaurante, do distrito e do entregador no momento da previsão.
Todas as variáveis usam apenas eventos anteriores ao momento da previsão. O nome e o cálculo de cada
variável estão nas Tabelas~\ref{tab:dic-orig} e \ref{tab:dic-var} (Apêndice~\ref{sec:dic}).

%=====================================================================
\section{Registro dos experimentos}
\label{sec:exp}

\paragraph{Gestão do ciclo de vida} Enquanto treinar um modelo se resume à execução do
código e obtenção de número, gerenciar o ciclo de vida implica conseguir responder depois, para
qualquer resultado: com quais dados, qual código e quais parâmetros ele foi gerado, quando, e
onde está o modelo. Com a ferramenta MLflow, foi possível guardar todos esses dados sobre cada
experimento desenvolvido: parâmetros, métricas, modelo e artefatos (gráficos de erro e importância das variáveis), junto com o hash dos
dados e o commit do código.

\paragraph{Métricas.} Como trata-se de um esperimento de regressão, usamos o \textbf{MAE} como métrica principal
pela sua interpretabilidade simples (média de minutos do erro). Além desta, também
mensuramos o \textbf{RMSE}, que mostra se há erros grandes,
que podem gerar reclamações, e o \textbf{R²}, que mostra quanto da variação do tempo o modelo
explica. Todas as métricas foram comparadas com a estimativa atual do aplicativo.

\subsection{Rodada 1 — pipeline com DVC e MLflow}

O treino foi realizado com 350 mil linhas e semente aleatória fixa (42), para que o resultado seja reprodutível.

\begin{table}[H]
\centering\footnotesize
\begin{tabular}{l l l l r r r}
\toprule
Execução & Dados & Modelo & Parâmetros principais & MAE & RMSE & R² \\
\midrule
01 & d1b18963 & Random Forest (9 var.) & 100 árvores, prof. 16 & 6,081 & 7,824 & 0,557 \\
02 & d1b18963 & Random Forest (19 var.) & 100 árvores, prof. 18 & 5,603 & 7,206 & 0,624 \\
03 & d1b18963 & Hist. Gradient Boosting (9) & 500 rodadas, taxa 0,08 & 5,941 & 7,654 & 0,576 \\
04 & d1b18963 & Hist. Gradient Boosting (19) & 500 rodadas, taxa 0,08 & 5,464 & 7,040 & 0,641 \\
05 & d1b18963 & LightGBM (9 var.) & 300 árvores, taxa 0,05 & 5,976 & 7,698 & 0,571 \\
06 & d1b18963 & LightGBM (19 var.) & 350 árvores, taxa 0,05 & 5,486 & 7,064 & 0,639 \\
07 & 127078ad & Estimativa do app & -- & 9,506 & 11,610 & 0,024 \\
\bottomrule
\end{tabular}
\caption{Execuções registradas no MLflow em 01/11/2026 (experimento \texttt{meituan-eta}),
avaliadas em \num{50000} pedidos de teste. ``Dados'' é o início do hash do arquivo usado;
9 variáveis no momento da compra, 19 no despacho. Fonte: \texttt{reports/experiment\_registry.csv}.}
\label{tab:r1}
\end{table}

\noindent\textbf{Problema encontrado:} nesta rodada, o tempo médio de cada restaurante foi
calculado usando todos os dados, antes de separar treino e teste. Assim, o conjunto de teste
teve acesso indireto a parte da resposta e os números ficaram melhores do que deveriam
(vazamento de dados). Mantivemos
esses números para comparar com as etapas anteriores, mas marcamos todas as execuções com o
aviso \texttt{known\_limitation}. O problema foi encontrado devido ao registro detalhado.

\subsection{Rodada 2 — consertado o vazamento, incluídas novas variáveis e treinado TabPFN}

Corrigimos o carregador de dados (\texttt{notebooks/dataloader.py}) para que passasse a calcular as médias
por restaurante e por entregador \emph{depois} de separar os dados, usando só o treino. Comparamos três tipos de modelo:
\begin{itemize}
  \item \textbf{Random Forest:} várias árvores de decisão independentes; a previsão é a média
        delas. Simples e estável, mas gera arquivos grandes.
  \item \textbf{LightGBM:} árvores construídas uma depois da outra; cada nova árvore corrige o
        erro das anteriores. Treina rapidamente em CPU e gera arquivos pequenos.
  \item \textbf{TabPFN-3.5:} uma rede neural que já vem treinada em milhões de tabelas
        artificiais. Precisa de GPU e fica bem mais lenta quando recebe mais exemplos.
\end{itemize}

\begin{table}[H]
\centering\footnotesize
\begin{tabular}{l r r r r r r}
\toprule
Modelo & Linhas & Var. & MAE & RMSE & R² & Tempo \\
\midrule
Estimativa do app Meituan & -- & -- & 9,535 & 11,646 & 0,029 & -- \\
Random Forest (50 árvores, prof. 12) & \num{377104} & 25 & 5,672 & 7,325 & 0,616 & 4 min \\
LightGBM & \num{377104} & 25 & 5,313 & 6,879 & 0,661 & 1 min \\
TabPFN-3.5 & \num{50000} & 25 & 5,323 & 6,962 & 0,653 & 7 min$^\dagger$ \\
LightGBM & \num{377104} & 38 & 5,133 & 6,652 & 0,683 & 2 min \\
TabPFN-3.5 & \num{100000} & 38 & 5,073 & 6,642 & 0,684 & 23 min$^\dagger$ \\
TabPFN-3.5 & \num{200000} & 38 & 5,006 & n/d & 0,692 & 106 min$^\dagger$ \\
\textbf{Combinação 0,80 TabPFN + 0,20 LightGBM} & \num{377104} & 38 & \textbf{5,001} & n/d & \textbf{0,694} & 108 min$^\dagger$ \\
\bottomrule
\end{tabular}
\caption{Rodada 2 treianando para prever tempo de entrega no momento do despacho com \num{53872} pedidos de teste. MAE e RMSE em minutos.
$^\dagger$GPU Tesla T4 no Google Colab; os demais foram executados em CPU.
n/d: o RMSE não ficou registrado nessa execução.}
\label{tab:r2}
\end{table}

\noindent A \emph{combinação} faz uma média ponderada das previsões dos dois modelos (80\% do
TabPFN e 20\% do LightGBM). Os pesos foram escolhidos na etapa de validação.

\paragraph{Versões registradas na rodada 2.} Pacote \texttt{tabpfn} 10.2.0 (modelo v3.5),
semente 42, divisão 70/20/10. Para cada cenário, o notebook salvou um manifesto com o hash
dos dados de treino (\texttt{72617a3a\ldots}) e de teste (\texttt{51e4e84d\ldots}), a
configuração e a GPU usada, além das previsões e de um arquivo com os resultados.

\begin{figure}[H]
\centering
\begin{tikzpicture}
\begin{axis}[
  width=0.62\linewidth, height=5.2cm,
  xmode=log, log basis x=10,
  xlabel={linhas usadas no treino}, ylabel={MAE no teste (min)},
  xtick={10000,25000,50000,100000}, xticklabels={10 mil,25 mil,50 mil,100 mil},
  ymin=5.0, ymax=5.6, ytick={5.0,5.1,5.2,5.3,5.4,5.5,5.6},
  yticklabel style={/pgf/number format/.cd, fixed, precision=1, use comma},
  grid=major, grid style={gray!20}, axis line style={gray!60}, tick style={gray!60},
  legend pos=north east, legend style={font=\footnotesize, draw=none, fill=none},
  every axis plot/.append style={line width=1.4pt, mark size=2.2pt}]
\addplot[color=corange, mark=*] coordinates {(10000,5.509) (25000,5.397) (50000,5.305) (75000,5.282) (100000,5.257)};
\addlegendentry{LightGBM}
\addplot[color=cblue, mark=*] coordinates {(10000,5.319) (25000,5.232) (50000,5.156) (75000,5.107) (100000,5.089)};
\addlegendentry{TabPFN-3.5}
\end{axis}
\end{tikzpicture}
\caption{Curva de aprendizado: os dois modelos recebem exatamente as mesmas linhas. Com a mesma
quantidade de dados, o TabPFN erra de 0,15 a 0,19 min a menos.}
\label{fig:lc}
\end{figure}

\paragraph{O que os testes mostraram.}
\begin{enumerate}
  \item As variáveis calculadas no despacho melhoram o desempenho dos modelos. No LightGBM, o erro caiu de 5,313 para 5,133 min.
  \item O TabPFN erra menos que o LightGBM mesmo usando só metade dos dados de treino. Cada vez que
        recebe o dobro de exemplos, o erro cai cerca de 0,07 min.
  \item A vantagem do TabPFN é pequena (0,13 min, cerca de 8 segundos por pedido) e foi
        medida numa única divisão dos dados.
  \item O custo é muito diferente: 106 minutos em GPU contra 2 minutos em CPU.
  \item O Random Forest errou mais nas duas rodadas e gerou artefatos de 176 a 330 MB.
\end{enumerate}

\noindent\textbf{O que ainda falta:} as execuções da rodada 2 foram feitas no Google Colab e
registradas nos manifestos, mas ainda não foram copiadas para o MLflow (Seção~\ref{sec:lim}).

\subsection{Escolha do modelo}

Para escolher o modelo a ser usado em produção, não basta ter o melhor desempenho.
O modelo utilizado será monitorado e tende a necessitar retreinamento devido a 
circunstâncias como \emph{data drift}. Por isso, levamos em consideração o desempenho, o custo, a facilidade
de treinar de novo, a interpretabilidade e a dependência de serviços externos.

\begin{table}[H]
\centering\small
\begin{tabularx}{\linewidth}{l c c c L}
\toprule
Papel & MAE & Treino & Arquivo & Motivo \\
\midrule
\textbf{Candidato: LightGBM} & 5,133 & 2 min, CPU & $\sim$1 MB & Desempenho próximo do melhor, baixo custo, indica o peso de cada variável, pode ser treinado novamente toda semana \\
Desafiante: TabPFN / combinação & 5,001 & 106 min, GPU & modelo externo & Menor erro; só substitui o candidato se o ganho se repetir em novos testes e o custo for aceito \\
Referência: Random Forest & 5,672 & 4 min, CPU & 176--330 MB & Modelo da demonstração do MLflow; serve de comparação \\
Contingência: estimativa do app & 9,535 & -- & -- & Usada se o modelo falhar ou for suspenso \\
\bottomrule
\end{tabularx}
\caption{Papel de cada modelo.}
\end{table}

%=====================================================================
\section{Model card}
\label{sec:card}

O model card é o documento que descreve o modelo para fins de uso, avaliação e auditoria,
contendo finalidade, dados de origem, desempenho, limites de validade, usos permitidos e
responsáveis. A versão completa (1.1) está no arquivo \texttt{MODEL\_CARD.md}, na pasta
\texttt{docs/model\_cards}, e contém também os hiperparâmetros, os identificadores de cada
execução no MLflow e as instruções de reprodução. Esta seção apresenta os dez itens exigidos
pela Etapa~4 do trabalho.

\subsection{Identificação}

\begin{table}[H]
\centering\small
\begin{tabularx}{\linewidth}{l L}
\toprule
Modelo & LightGBM (regressão), 38 variáveis, calculado no momento do despacho; versão equivalente com 13 variáveis para o momento da compra \\
Versão do card & 1.1, de 22/11/2026 \\
Status & experimental; não aprovado para uso em produção \\
Modelos relacionados & TabPFN-3.5 e combinação TabPFN + LightGBM (desafiantes); Random Forest (referência) \\
\bottomrule
\end{tabularx}
\end{table}

\subsection{Itens do Model Card}

\paragraph{1. Objetivo.} Estimar, em minutos, o tempo entre a criação do pedido e a chegada ao
cliente, em dois momentos: na compra (antes de haver entregador) e no despacho (após o aceite
do entregador). A estimativa apoia o planejamento da operação e a comunicação com o cliente.

\paragraph{2. Dados utilizados.} Base pública da Meituan (INFORMS TSL Challenge), 8 dias de
outubro de 2022, \num{538435} entregas válidas após a limpeza descrita na
Seção~\ref{sec:dados}. Licença CC BY-NC 4.0. Treino com \num{377104} linhas, validação com
\num{107744} e teste com \num{53872}.

\paragraph{3. População e contexto.} Pedidos imediatos, aceitos por um entregador e
concluídos, com duração entre 5 e 90 minutos, em uma única cidade chinesa. Ficam fora pedidos
agendados, cancelados ou recusados por todos os entregadores. O modelo descreve esse
contexto específico e não deve ser tratado como representativo de outras cidades, outros
períodos do ano ou outras empresas.

\paragraph{4. Métricas.} MAE (métrica principal), RMSE e R² comparados com a estimativa
atual do aplicativo no mesmo conjunto de teste.

\begin{table}[H]
\centering\small
\begin{tabular}{l r r r}
\toprule
& MAE (min) & RMSE (min) & R² \\
\midrule
LightGBM, despacho (38 variáveis) & 5,133 & 6,652 & 0,683 \\
Estimativa do aplicativo & 9,535 & 11,646 & 0,029 \\
Diferença & $-46\%$ & $-43\%$ & $+0{,}654$ \\
\bottomrule
\end{tabular}
\end{table}

\paragraph{5. Limitações.}
\begin{itemize}
  \item Apenas 8 dias e uma cidade; não há validação em outro período ou região real.
  \item Uma única divisão aleatória dos dados, sem repetição com outras sementes.
  \item Ausência de informações sobre clima, trânsito e promoções.
  \item O modelo fornece um valor único, sem intervalo de confiança; além disso, não há garantia de que o pedido
        chegue antes do tempo previsto.
  \item As coordenadas foram deslocadas pela Meituan; as distâncias são aproximações.
  \item Na rodada 1, o vazamento de dados torna as métricas otimistas; a rodada 2 corrige o
        problema, mas suas execuções ainda não foram migradas para o MLflow.
\end{itemize}

\paragraph{6. Usos recomendados.} Apoio a analistas de operação; planejamento da quantidade de
entregadores por região e horário; comparação entre regiões da cidade; base para o
desenvolvimento de um modelo de produção; ensino de práticas de reprodutibilidade. Em todos
os casos, com supervisão humana.

\paragraph{7. Usos não recomendados.}
\begin{itemize}
  \item \textbf{Uso em operação sem um período de validação prévio:} antes de afetar clientes,
        o modelo deve rodar em paralelo com dados atuais, e o erro observado deve ser comparado
        com o do teste.
  \item \textbf{Uso comercial sem ampliar a base de treino:} 8 dias de uma única cidade não
        representam a variação de uma operação real (estações do ano, feriados, outras
        regiões). Além disso, a licença CC BY-NC dos dados atuais não permite uso comercial.
  \item \textbf{Avaliar, punir ou ranquear entregadores:} o tempo de entrega depende de fatores
        fora do controle do entregador (região, restaurante, demanda), portanto entregadores que
        levem mais tempo na entrega do que o previsto pelo modelo não podem ser punidos por isso.
\end{itemize}

\paragraph{8. Riscos.} Prometer prazos curtos demais e gerar atrasos; perda de qualidade
quando os dados mudam (data drift); erro maior em algumas regiões; confiança excessiva nas
previsões. A matriz completa, com probabilidade, impacto, ação e responsável, está na
Seção~\ref{sec:risco}.

\paragraph{9. Aspectos éticos.}
\begin{itemize}
  \item \textbf{Privacidade:} os dados contêm localização de clientes, restaurantes e
        entregadores. Mesmo com coordenadas deslocadas, há risco de reidentificação; os dados
        devem ficar fora do repositório e com acesso restrito.
  \item \textbf{Entregadores:} as variáveis com o histórico de cada entregador melhoram a
        previsão, mas não podem ser usadas para avaliação individual.
  \item \textbf{Equidade:} o modelo não usa atributos pessoais, mas região e restaurante podem
        refletir diferenças socioeconômicas. O erro por distrito faz parte do monitoramento.
  \item \textbf{Transparência e revisão humana:} limitações documentadas; decisões com efeito
        sobre clientes ou entregadores exigem revisão humana.
\end{itemize}

\paragraph{10. Responsáveis.} Equipe do projeto: Guilherme Braga, Guilherme Costanzi e Rafael Cobo. A aprovação para uso cabe ao comitê de governança
(Seção~\ref{sec:gov}). O card deve ser revisado a cada novo treino, após incidentes e, no
mínimo, a cada seis meses.

%=====================================================================
\section{Plano de monitoramento}
\label{sec:mon}

Para garantir que a manutenção do funcionamento do modelo após ele entrar em uso, planejamos medir
quatro grupos de indicadores toda semana (Tabelas~\ref{tab:mon} e \ref{tab:mon2}). Os gatilhos de ação 
foram definidos a partir dos resultados do LightGBM candidato (MAE 5,13 min; R² 0,683). 
O tempo real de cada entrega é determinado logo depois do pedido, então o erro do modelo pode ser medido imediatamente.

O data drift será mensurado a partir do índice de estabilidade populacional (PSI), que detecta mudanças
na distribuição dos valores de um atributo em um conjunto de dados.

\begin{table}[H]
\centering\footnotesize
\begin{tabularx}{\linewidth}{l >{\raggedright\arraybackslash}p{4.2cm} L L}
\toprule
Grupo & Indicador & Gatilho & Ação \\
\midrule
Dados & Valores faltantes ou zeros & $>1\%$ numa variável ou o dobro do normal & Investigar a origem e usar a estimativa do app (contingência) se for uma variável importante \\
Dados & Data drift: mudança nas variáveis de entrada (índice PSI) & PSI $>0{,}10$ alerta; $>0{,}25$ crítico & Ver quais variáveis mudaram; no crítico, avaliar novo treino \\
Dados & Mudança na distribuição do tempo de entrega real (média e variação) & mudança > 10\% & Conferir junto com o R², investigar a causa \\
Dados & Qualidade e formato & coluna faltando, tipo errado, hora fora de 0--23, distância negativa & Bloquear o lote e avisar o responsável de dados \\
Dados & Categorias novas & $>5\%$ dos pedidos com distrito ou restaurante desconhecido & Incluir no próximo treino \\
\midrule
Modelo & Desempenho: MAE semanal & $>5{,}9$ min ($+15\%$) alerta; $>6{,}4$ min ($+25\%$) crítico & Alerta: analisar; crítico: treinar de novo e reavaliar \\
Modelo & Erros: desvio médio (real $-$ previsto) & mais de 1 min para um dos lados & Ajustar o modelo; ver se a operação mudou \\
Modelo & Erros grandes & pedidos com erro $>10$ min acima de 1,5 vez o normal & Revisar os casos por região e horário \\
Modelo & Model drift: queda do R² & queda $>0{,}05$ alerta; $>0{,}15$ crítico & Comparar com o desafiante; treinar de novo \\
Modelo & Erro por região & diferença de MAE $>2$ min entre distritos & Análise de viés pelo comitê \\
\bottomrule
\end{tabularx}
\caption{Plano de monitoramento (parte 1): dados e modelo.}
\label{tab:mon}
\end{table}

\begin{table}[H]
\centering\footnotesize
\begin{tabularx}{\linewidth}{l >{\raggedright\arraybackslash}p{4.2cm} L L}
\toprule
Grupo & Indicador & Limite & Ação \\
\midrule
Negócio & Benefício esperado: ganho sobre a estimativa do app & MAE menos de 40\% menor que o do app & Reavaliar se o uso do modelo se justifica \\
Negócio & Pedidos que chegam depois do prazo prometido & $>17\%$ (hoje $\approx 14\%$) & Rever a folga do prazo prometido \\
Negócio & Reclamações por atraso a cada 1.000 pedidos & $+20\%$ sobre a média de 4 semanas & Analisar a causa com a operação \\
Negócio & Pedidos com previsão e tempo de resposta & $<99\%$ dos pedidos; resposta lenta ($>200$ ms) & Verificar o serviço; usar a estimativa do app (contingência) \\
\midrule
Governança & Incidentes & Qualquer incidente grave & Registrar, suspender o uso automático, analisar em 5 dias úteis \\
Governança & Reclamações sobre o modelo & Qualquer reclamação de cliente ou entregador ligada à previsão & Analisar e responder com registro \\
Governança & Vieses & Diferença por região ou turno acima do limite & Rever variáveis e regras de uso \\
Governança & Necessidade de revisão & Model card com mais de 6 meses, novo treino ou incidente & Revisão obrigatória antes de continuar \\
\bottomrule
\end{tabularx}
\caption{Plano de monitoramento (parte 2): negócio e governança.}
\label{tab:mon2}
\end{table}

\subsection{Simulação de monitoramento com a chegada de dados de uma região nova}

Não temos dados de outras semanas. Por isso, simulamos a chegada de dados novos 
treinando o modelo sem os distritos 8 e 9 e testando só com eles, como se houvesse uma expansão
para uma região que o modelo nunca viu. Usamos o TabPFN com 50 mil linhas e 25 variáveis, que já
tinha sido rodado nos dois cenários.

\begin{table}[H]
\centering\small
\begin{tabular}{l r r l}
\toprule
Indicador & Referência & Região nova & Situação \\
\midrule
MAE (min) & 5,323 & 5,706 ($+7{,}2\%$) & dentro do limite \\
R² & 0,653 & 0,548 ($-0{,}105$) & \textbf{alerta} \\
Variação do tempo de entrega real (min²) & 139,7 & 120,9 ($-13\%$) & \textbf{mudou} \\
R² da estimativa do app & 0,029 & $-0{,}117$ & piora também \\
\bottomrule
\end{tabular}
\caption{Simulação de chegada de dados novos. Referência: divisão aleatória (\num{53872}
pedidos); região nova: distritos 8 e 9 (\num{79023} pedidos).}
\label{tab:sim}
\end{table}

\noindent\textbf{Interpretação.} Considerando apenas o MAE, nenhum alerta seria emitido, mas o plano
de monitoramento identifica a mudança na distribuição do tempo de entrega: na região
nova os tempos variam menos, e o modelo explica uma parte menor dessa variação. A ação prevista
é abrir uma revisão, analisar o erro por distrito e incluir dados da região no próximo treino
antes de usar o modelo nela. O exemplo mostra por que o plano combina vários indicadores.

%=====================================================================
\section{Matriz de riscos}
\label{sec:risco}

Cada risco mapeado recebe uma nota de 1 (baixa) a 3 (alta) para a probabilidade (P) de acontecer e
para o impacto (I) se acontecer. O nível é P $\times$ I: 6 ou mais é alto, 3 ou 4 é médio,
até 2 é baixo.

\begin{table}[H]
\centering\footnotesize
\begin{tabularx}{\linewidth}{l L c c c L l}
\toprule
Tipo & Risco & P & I & Nível & O que fazemos & Dono \\
\midrule
Dados & a demanda, os entregadores ou a cidade mudam (data drift) & Alta & Alto & \Ra & PSI semanal, novo treino quando passar do limite, simulação por distrito & Dados \\
Dados & amostra só com pedidos concluídos (faltam os cancelados) & Alta & Médio & \Ra & avisar no model card; incluir cancelados quando houver dados & Dados \\
Dados & zeros no lugar de dados faltantes e erros de formato & Média & Médio & \Rm & etapa de validação; bloqueio do lote & Dados \\
Modelo & subestimar o tempo e prometer um prazo curto demais & Média & Alto & \Ra & não usar como prazo direto; somar folga ou prever uma faixa & Modelo \\
Modelo & resultados melhores do que a realidade (vazamento, uma só divisão) & Média & Médio & \Rm & médias calculadas só com o treino (rodada 2); repetir os testes & Modelo \\
Modelo & dependência de GPU e de pacote externo (TabPFN) & Média & Médio & \Rm & LightGBM como candidato; TabPFN só como desafiante; versão fixa & Modelo \\
Privacidade & reidentificação de pessoas pela localização & Baixa & Alto & \Rm & dados fora do Git; acesso restrito; não divulgar coordenadas & Dados \\
Privacidade & uso fora da licença (comercial) & Baixa & Alto & \Rm & só uso acadêmico; citar a Meituan & Equipe \\
Segurança & arquivo do modelo (\texttt{.pkl}) alterado por terceiros & Baixa & Alto & \Rm & só carregar arquivos com o hash registrado no \texttt{dvc.lock} & Modelo \\
Segurança & senha de acesso (token) do TabPFN exposta & Média & Médio & \Rm & guardar em local protegido, nunca no notebook ou no Git & Equipe \\
Ética & histórico do entregador usado para avaliá-lo ou puni-lo & Média & Alto & \Ra & proibido no model card; nunca criar ranking; testar o modelo sem essas variáveis & Comitê \\
Ética & erro maior em algumas regiões (serviço pior) & Média & Médio & \Rm & erro por distrito no monitoramento; análise de viés & Comitê \\
Ética & confiança excessiva nas previsões do modelo & Média & Médio & \Rm & mostrar a margem de erro; decisões importantes com revisão humana & Comitê \\
\bottomrule
\end{tabularx}
\caption{Matriz de riscos de dados, modelo, privacidade, segurança e ética.}
\label{tab:risco}
\end{table}

%=====================================================================
\section{Como a solução será governada}
\label{sec:gov}

\paragraph{Responsabilidades previstas pela solução}
\begin{itemize}
  \item \textbf{Responsável de dados:} qualidade, validação, acesso e licença dos dados.
  \item \textbf{Responsável do modelo:} treino, avaliação, MLflow, model card e acompanhamento
        dos indicadores.
  \item \textbf{Comitê de governança} (produto, operação e jurídico/privacidade): aprova cada
        versão, analisa incidentes e vieses e decide os usos permitidos.
\end{itemize}
Nessa disciplina, os três papéis foram exercidos pelo grupo, mas numa empresa seriam
exercidos por pessoas diferentes.

\paragraph{Protocolo para implementação de uma nova versão do modelo}
\begin{enumerate}
  \item a mudança no código, nos dados ou nos parâmetros é registrada no Git e no DVC;
  \item o \texttt{dvc repro} gera o novo modelo e o MLflow registra a execução;
  \item a nova versão é comparada com a atual, nos mesmos dados de teste e por região;
  \item o model card é atualizado e o comitê aprova ou recusa, com registro da decisão;
  \item a versão aprovada roda primeiro ``em paralelo'', fazendo previsões mas sem mostrar elas
        para o cliente, para fins de monitoramento e avaliação em ambiente de produção;
  \item o plano de monitoramento (Tabela~\ref{tab:mon}) acompanha a versão; num alerta
        crítico, volta-se ao passo 1 ou à versão anterior, identificada no MLflow.
\end{enumerate}

%=====================================================================
\section{Limitações}
\label{sec:lim}

\begin{itemize}
  \item Os dados cobrem só 8 dias e uma cidade; o modelo não foi testado em outro período ou
        região real.
  \item Os resultados vêm de uma única divisão aleatória dos dados; diferenças pequenas entre
        modelos (como os 0,13 min entre TabPFN e LightGBM) podem mudar com outra divisão.
  \item Não há clima, trânsito ou promoções nos dados.
  \item O modelo dá um valor único, sem margem de erro.
  \item As execuções da rodada 2 ainda não foram migradas para o MLflow.
  \item O plano de monitoramento foi definido e simulado, mas não está funcionando de forma
        automática.
\end{itemize}

%=====================================================================
\section{Onde está cada item da entrega final}
\label{sec:checklist}

\begin{table}[H]
\centering\small
\begin{tabularx}{\linewidth}{l L}
\toprule
Item pedido & Onde está \\
\midrule
1. Repositório da solução & Seção~\ref{sec:repo}; \texttt{README.md}, \texttt{dvc.yaml}, \texttt{meituan\_eta/}, \texttt{notebooks/} \\
2. Dataset documentado & Seção~\ref{sec:dados} e Apêndice~\ref{sec:dic}; \texttt{quality/}; relatório da Etapa 2 \\
3. Registro dos experimentos & Seção~\ref{sec:exp}; MLflow; \texttt{reports/experiment\_registry.csv} \\
4. Model card & Seção~\ref{sec:card}; \texttt{docs/model\_cards/MODEL\_CARD.md} \\
5. Plano de monitoramento & Seção~\ref{sec:mon} (Tabelas~\ref{tab:mon} e \ref{tab:mon2} e simulação) \\
6. Matriz de riscos & Seção~\ref{sec:risco} \\
7. Relatório executivo & Seções~\ref{sec:exec} e \ref{sec:gov} \\
\bottomrule
\end{tabularx}
\end{table}

\vspace{0.5em}
\noindent{\small\emph{This research was supported by data provided by Meituan.}}

%=====================================================================
\appendix
\section{Dicionário de dados}
\label{sec:dic}

\begin{table}[H]
\centering\footnotesize
\begin{tabularx}{\linewidth}{l l L}
\toprule
Coluna original & Tipo & Significado \\
\midrule
\texttt{dt} & data & dia do registro (AAAAMMDD) \\
\texttt{order\_id}, \texttt{waybill\_id} & código & pedido e tentativa de entrega \\
\texttt{courier\_id}, \texttt{poi\_id}, \texttt{da\_id} & código & entregador, restaurante e distrito \\
\texttt{is\_courier\_grabbed} & 0 ou 1 & o entregador aceitou o pedido \\
\texttt{is\_weekend}, \texttt{is\_prebook} & 0 ou 1 & fim de semana; pedido agendado \\
\texttt{sender\_lat/lng} & coordenada & localização do restaurante (deslocada) \\
\texttt{recipient\_lat/lng} & coordenada & localização do cliente (deslocada) \\
\texttt{grab\_lat/lng} & coordenada & onde o entregador estava ao aceitar (deslocada) \\
\texttt{platform\_order\_time} & horário & criação do pedido \\
\texttt{order\_push\_time}, \texttt{dispatch\_time} & horário & envio ao sistema; oferta ao entregador \\
\texttt{grab\_time}, \texttt{fetch\_time}, \texttt{arrive\_time} & horário & aceite; retirada no restaurante; chegada ao cliente \\
\texttt{estimate\_meal\_prepare\_time} & horário & quando a cozinha estima que o prato fica pronto \\
\texttt{estimate\_arrived\_time} & horário & prazo prometido ao cliente (a estimativa do app) \\
\bottomrule
\end{tabularx}
\caption{As 23 colunas da tabela original de entregas. Horários em segundos desde 1970 (Unix).}
\label{tab:dic-orig}
\end{table}

\begin{table}[H]
\centering\scriptsize
\begin{tabularx}{\linewidth}{>{\raggedright\arraybackslash}p{5.6cm} l L}
\toprule
Variável & Momento & Como é calculada \\
\midrule
\texttt{target\_total\_time\_min} & alvo & (\texttt{arrive\_time} $-$ \texttt{platform\_order\_time}) / 60 \\
\texttt{meituan\_baseline\_min} & referência & (\texttt{estimate\_arrived\_time} $-$ \texttt{platform\_order\_time}) / 60 \\
\midrule
\multicolumn{3}{l}{\emph{Compra: \texttt{FEATURES\_REDUCED} (13)}} \\
\texttt{dist\_rest\_client\_km} & compra & distância de Manhattan restaurante $\rightarrow$ cliente \\
\texttt{restaurant\_lat}, \texttt{restaurant\_lng} & compra & coordenadas do restaurante (deslocadas) \\
\texttt{customer\_lat}, \texttt{customer\_lng} & compra & coordenadas do cliente (deslocadas) \\
\texttt{order\_hour}, \texttt{day\_of\_week}, \texttt{is\_weekend} & compra & hora, dia da semana e fim de semana do pedido (UTC+8) \\
\texttt{sin\_order\_hour}, \texttt{cos\_order\_hour} & compra & minuto do dia em seno e cosseno \\
\texttt{est\_meal\_prep\_min} & compra & (\texttt{estimate\_meal\_prepare\_time} $-$ criação do pedido) / 60 \\
\texttt{district\_id} & compra & código do distrito (categórica) \\
\texttt{restaurant\_target\_enc} & compra & tempo médio do restaurante, suavizado, calculado só com o treino \\
\midrule
\multicolumn{3}{l}{\emph{Despacho: \texttt{FEATURES\_COMPLETE} = compra + 12 (25)}} \\
\texttt{courier\_lat}, \texttt{courier\_lng} & despacho & posição do entregador no aceite (deslocada) \\
\texttt{dist\_driver\_rest\_km} & despacho & distância de Manhattan entregador $\rightarrow$ restaurante \\
\texttt{dist\_total\_driver\_trip\_km} & despacho & entregador $\rightarrow$ restaurante $\rightarrow$ cliente \\
\texttt{time\_to\_grab\_min} & despacho & (\texttt{grab\_time} $-$ criação do pedido) / 60 \\
\texttt{queue\_wait\_min} & despacho & (\texttt{dispatch\_time} $-$ criação do pedido) / 60 \\
\texttt{rem\_prep\_at\_dispatch\_min} & despacho & (\texttt{estimate\_meal\_prepare\_time} $-$ \texttt{dispatch\_time}) / 60 \\
\texttt{prior\_rejections} & despacho & recusas anteriores do mesmo pedido \\
\texttt{courier\_active\_orders} & despacho & pedidos em andamento com o entregador no aceite \\
\texttt{district\_demand\_20min} & despacho & pedidos aceitos no distrito nos 20 min anteriores \\
\texttt{is\_in\_wave}, \texttt{wave\_position} & despacho & se o pedido faz parte de uma onda de entregas e sua posição nela \\
\midrule
\multicolumn{3}{l}{\emph{Extras da rodada 2: \texttt{FEATURES\_COMPLETE\_EXTRA} = despacho + 13 (38)}} \\
\texttt{restaurant\_active\_orders} & compra & pedidos do restaurante criados e ainda não retirados \\
\texttt{district\_recent\_eta\_30min} & compra & tempo médio das entregas concluídas no distrito nos 30 min anteriores \\
\texttt{district\_recent\_delay\_30min} & compra & atraso médio (real $-$ prometido) no distrito nos 30 min anteriores \\
\texttt{district\_demand\_order\_20min} & compra & pedidos criados no distrito nos 20 min anteriores \\
\texttt{district\_active\_couriers\_order\_20min} & compra & entregadores distintos ativos no distrito nos 20 min anteriores \\
\texttt{supply\_demand\_ratio\_order} & compra & pedidos / (entregadores ativos + 1) \\
\texttt{restaurant\_recent\_prep\_min\_60min} & compra & tempo médio de preparo do restaurante na última hora \\
\texttt{district\_active\_couriers\_20min} & despacho & entregadores ativos no distrito nos 20 min antes do aceite \\
\texttt{supply\_demand\_ratio} & despacho & demanda / (entregadores ativos + 1) no aceite \\
\texttt{district\_recent\_eta\_dispatch\_30min} & despacho & tempo médio recente no distrito, medido no aceite \\
\texttt{district\_recent\_delay\_dispatch\_30min} & despacho & atraso médio recente no distrito, medido no aceite \\
\texttt{courier\_target\_enc} & despacho & tempo médio histórico do entregador, suavizado, só com o treino \\
\texttt{courier\_speed\_enc} & despacho & velocidade média histórica do entregador (km/h), só com o treino \\
\bottomrule
\end{tabularx}
\caption{Variáveis criadas, com os nomes usados em \texttt{notebooks/dataloader.py}. Na rodada 1
(pipeline \texttt{meituan\_eta}), algumas têm nomes em português (por exemplo,
\texttt{tempo\_ate\_aceite\_min}, \texttt{carga\_entregador}, \texttt{poi\_target\_enc}).}
\label{tab:dic-var}
\end{table}

\end{document}
